\chapter{Dove gira il sistema}
\label{cap:deployment}

\section{Ambiente A -- sul nostro computer}

Usiamo \textbf{Kubernetes}, un programma che avvia i servizi, li tiene d'occhio e li riavvia da solo quando si rompono. Lo facciamo girare sul computer con \textbf{kind}, che simula tre macchine dentro Docker:

\begin{description}
  \item[control-plane] il ``cervello'' di Kubernetes e la porta d'ingresso; non lo rompiamo mai.
  \item[worker-1] ospita database e coda.
  \item[worker-2] ospita le altre copie; è la macchina che spegniamo in uno dei test.
\end{description}

\begin{figure}[htbp]
\centering
\resizebox{0.9\linewidth}{!}{%
\begin{tikzpicture}[node distance=0.35cm]
\node[box] (pg) {Database};
\node[box, below=of pg] (mq) {Coda};
\node[box, below=of mq] (a1) {API, worker};
\node[draw=unina, dashed, rounded corners, fit=(pg)(mq)(a1), inner sep=8pt,
      label={[font=\small\color{unina}]above:worker-1}] (w1) {};
\node[box, right=2.2cm of pg] (a2) {API, worker};
\node[draw=unina, dashed, rounded corners, fit=(a2), inner sep=8pt,
      label={[font=\small\color{unina}]above:worker-2}] (w2) {};
\node[box, left=2.4cm of pg] (cp) {Cervello di\\Kubernetes};
\node[draw=unina, dashed, rounded corners, fit=(cp), inner sep=8pt,
      label={[font=\small\color{unina}]above:control-plane}] (c) {};
\node[box, below=2.2cm of cp] (jm) {Generatore\\di ordini};
\draw[freccia] (jm) -- (c);
\draw[freccia] (c.east) -- (w1.west);
\end{tikzpicture}}
\caption{L'ambiente locale: tre macchine simulate.}
\label{fig:deploy-locale}
\end{figure}

Il file \texttt{deploy/k8s/kind-cluster.yaml} descrive queste tre macchine: chiunque può ricrearle identiche con un comando.

\paragraph{I due computer del gruppo.} Il Mac di F.~Terrecuso (8 GB) serve per sviluppare e provare. Le misure vere si fanno sul computer di G.~Vellotti (16 GB), così la memoria scarsa non falsa i risultati (decisione D0013).

\section{Ambiente B -- nel cloud (AWS)}

Ripetiamo una parte degli esperimenti su AWS, con tre macchine vere in tre edifici diversi. I servizi usati:

\begin{description}
  \item[EKS] il ``cervello'' di Kubernetes, gestito da Amazon.
  \item[EC2] le macchine: tre per il sistema, una per generare gli ordini.
  \item[ALB] il bilanciatore che smista le richieste.
  \item[EBS] i dischi del database e della coda.
\end{description}

Costo stimato dell'intera fase: circa 15--20 dollari, coperti dai crediti gratuiti di un account nuovo. Il rischio vero è dimenticare risorse accese: ogni sessione termina con una verifica che tutto sia stato spento.
